\documentclass[11pt,twoside,openany]{book} % "book" es mejor para manuales largos; "twoside" para formato libro.

% --- PAQUETES DE IDIOMA Y FORMATO ---
\usepackage[utf8]{inputenc}
\usepackage[spanish]{babel}
\usepackage[T1]{fontenc}
\usepackage{titlesec} % Para personalizar títulos de capítulos estilo 40k

% --- PAQUETES DE DISEÑO ---
\usepackage{geometry}
\geometry{top=2cm, bottom=2cm, left=2cm, right=2cm, headheight=14pt}
\usepackage{graphicx}
\usepackage{xcolor}
\usepackage{tcolorbox} % Imprescindible para las reglas especiales
\tcbuselibrary{skins, breakable}
\usepackage{fancyhdr} % Para encabezados con el nombre del manual
\usepackage{multicol} % Los manuales de 40k suelen ir a 2 columnas

% --- PAQUETES DE MATEMÁTICAS Y TABLAS ---
\usepackage{amsmath, amssymb}
\usepackage{tabularx}
\usepackage{booktabs}
\usepackage[table]{xcolor}

% --- LIBRERÍAS ADICIONALES PARA LA FICHA ---
\tcbuselibrary{skins, hooks}

% --- PLANTILLA DE HÉROE LEGENDARIO ---
\newtcolorbox{HeroSheet}[2]{
    enhanced,
    colback=white,
    colframe=DeepBlack,
    fonttitle=\bfseries\Large,
    title={#1 \hfill \small #2}, % Nombre del Héroe y Puntos/Rol
    attach boxed title to top left={xshift=5mm, yshift=-3mm},
    boxed title style={colback=WarhammerRed, sharp corners},
    sharp corners,
    drop shadow,
    before skip=10pt,
    after skip=10pt,
    underlay={ % Línea decorativa roja lateral
        \path[fill=WarhammerRed] ([xshift=1mm]frame.north west) -- ([xshift=2mm]frame.north west) -- ([xshift=2mm]frame.south west) -- ([xshift=1mm]frame.south west) -- cycle;
    }
}

% Comandos internos para organizar la información dentro de la ficha
\newcommand{\heroDescription}[1]{{\small\itshape #1}\par\medskip\hrule\medskip}
\newcommand{\abilitiesHeader}{\paragraph{\color{WarhammerRed}HABILIDADES Y PODERES}}


% --- DEFINICIÓN DE COLORES TEMÁTICOS ---
\definecolor{WarhammerRed}{HTML}{8E0B0B}
\definecolor{GoldRule}{HTML}{D4AF37}
\definecolor{DeepBlack}{HTML}{1A1A1A}

% --- CONFIGURACIÓN DE TÍTULOS (ESTILO D&D) ---
\titleformat{\chapter}[display]
  {\normalfont\huge\bfseries\color{WarhammerRed}}
  {\chaptertitlename\ \thechapter}{10pt}{\Huge}

% --- COMANDO PARA REGLAS ESPECIALES (CAJA ESTILO WH40K) ---
\definecolor{GrimRed}{HTML}{4A0000} % Rojo sangre oscuro
\definecolor{LanzaGrey}{HTML}{2F2F2F} % Gris metal gastado

% Modifica tu reglaespecial para que se vea más agresiva
\newtcolorbox{reglaespecial}[1]{
    colback=black!5, 
    colframe=GrimRed,
    fonttitle=\bfseries\color{white},
    title=#1,
    enhanced,
    sharp corners, % Las esquinas afiladas dan sensación de peligro
    borderline={0.5pt}{0pt}{GrimRed}
}

% --- COMANDO PARA DATASHEET (ESTILO ATRIBUTOS) ---
\newcommand{\stats}[6]{
    \begin{center}
    \begin{tabularx}{\linewidth}{|X|X|X|X|X|X|}
    \hline
    \rowcolor{gray!20} \textbf{MOV} & \textbf{MELEE} & \textbf{RANGO} & \textbf{REF} & \textbf{BLOQ} & \textbf{FOCUS} \\ \hline
    #1 & #2 & #3 & #4 & #5 & #6 \\ \hline
    \end{tabularx}
    \end{center}
}

% Estilo para Relatos de Héroes (Fluff)
\newtcolorbox{relato}[1]{
    colback=black!90, % Fondo casi negro
    colframe=GoldRule, % Borde dorado
    fonttitle=\bfseries\scshape,
    colupper=white,
    title=#1,
    enhanced,
    breakable,
    fontupper=\itshape, % Texto en cursiva para narrativa
    sharp corners
}

% --- INICIO DEL DOCUMENTO ---
\begin{document}

% 1. Carátula (La que ya tienes o una más artística)
% --- ARCHIVO: caratula.tex ---
\begin{titlepage}
    \centering
    
    % Fondo o Marco (Opcional: Si tienes una imagen de fondo, se descomenta lo siguiente)
    % \begin{tikzpicture}[remember picture,overlay]
    %     \node[anchor=center] at (current page.center) {\includegraphics[width=\paperwidth,height=\paperheight]{Imagenes/fondo_epico.jpg}};
    % \end{tikzpicture}

    % Espacio superior para que el título no quede pegado arriba
    \vspace*{2cm}

    % Línea decorativa superior estilo Manual
    \noindent\makebox[\linewidth]{\rule{\textwidth}{2pt}} 
    \vspace{0.4cm}

    % TÍTULO PRINCIPAL
    {\Huge \bfseries \color{WarhammerRed} LA ERA DEL MARTILLO \par}
    
    \vspace{0.2cm}
    
    % SUBTÍTULO O VERSIÓN
    {\Large \scshape Guía Practica Para La Guerra \par}
    
    \vspace{0.4cm}
    \noindent\makebox[\linewidth]{\rule{\textwidth}{2pt}}

    \vspace{2cm}

    % IMAGEN CENTRAL (Aquí iría el logo de tu juego o un arte conceptual)
    % Si no tienes imagen aún, puedes dejar el espacio o usar un símbolo TikZ
\begin{tikzpicture}
    % Definición de Colores
    % \definecolor{GoldRule}{HTML}{D4AF37}
    % \definecolor{WarhammerRed}{HTML}{8E0B0B}
    % \definecolor{DeepBlack}{HTML}{1A1A1A}

    % 1. El Triángulo Invertido (La Punta de la Lanza)
    \draw[line width=4pt, color=GoldRule, line join=round] (0,3.46) -- (4,3.46) -- (2,0) -- cycle;

    % 2. El Gran Rayo Central (Canalización de Energía)
    \fill[color=WarhammerRed] 
        (1.2, 3.2)  -- (2.8, 3.2)   % Base superior
        -- (2.2, 2.1)               % Quiebre
        -- (2.6, 2.0)               % Saliente
        -- (1.9, 0.6)               % Vértice inferior del rayo
        -- (2.0, 1.6)               % Saliente
        -- (1.4, 1.7)               % Quiebre
        -- cycle;

    % 3. El Martillo Negro Flat (Debajo de la punta)
    % Dibujado de forma minimalista: mango y cabeza rectangular
    \begin{scope}[shift={(2, -0.3)}] % Posicionamiento debajo de la punta
        % Cabeza del Martillo
        \fill[color=DeepBlack] (-0.8, -0.4) rectangle (0.8, 0.1); 
        % Mango del Martillo
        \fill[color=DeepBlack] (-0.15, -1.6) rectangle (0.15, -0.4);
        % Detalle de pomo (opcional, para estilo)
        \fill[color=DeepBlack] (-0.2, -1.7) rectangle (0.2, -1.6);
    \end{scope}

\end{tikzpicture}

    \vfill

    % SECCIÓN INFERIOR: AUTORÍA Y CRÉDITOS
    \begin{tcolorbox}[colframe=WarhammerRed, colback=black!90, arc=0mm, width=0.8\textwidth]
        \centering
        {\color{white} \large \textbf{DISEÑADO POR:} \\ \scshape Franco Joaquín Gómez}
    \end{tcolorbox}

    \vspace{1cm}

    % FECHA O VERSIÓN DEL DOCUMENTO
    {\small \color{gray} Versión 0.0.1 - ''EL INICIO DE LA ERA'' \\ \the\year \par}

\end{titlepage}

\frontmatter
\tableofcontents

\mainmatter
\pagestyle{fancy}
\fancyhead[LE,RO]{\thepage}
\fancyhead[RE,LO]{LA ERA DEL MARTILLO - REGLAS CENTRALES}

% 2. Capítulos del Manual
\columnsep=0.7cm
\begin{multicols*}{2} % Activa las dos columnas estilo manual oficial

% --- ARCHIVO: Capitulos/introduccion.tex ---

\chapter{Sistema Lanza - Keitros}

\begin{quote}
\textit{''En este planeta llamado Keitros, extraños y desconocidos dioses poblaron estas tierras, y crearon seres de diferentes formas, cada dios se enfoco un grupo de ellos, y sin saber como o porque, cultivo a estos seres. Los Dioses les dieron poder e intelecto, sentir el poder. Y cuando estas creaturas se encontraron por primera vez en el mundo, ''La Era Del Martillo'' comenzó.''}
\end{quote}


\section{¿Qué es La Era del Martillo?}
\textbf{La Era del Martillo} es un juego de escaramuzas heroicas diseñado para jugadores que buscan profundidad táctica y consecuencias reales en cada decisión. A diferencia de otros juegos de guerra, aquí cada unidad es un activo invaluable. El sistema no se trata solo de mover miniaturas, sino de gestionar el cansancio, la posición y, sobre todo, la energía del entorno.

\begin{reglaespecial}{La Filosofía del Juego}
Este sistema se basa en tres pilares fundamentales:
\begin{itemize}
    \item \textbf{Letalidad Detallada:} Gracias al sistema de localización de daños, un solo golpe bien colocado puede cambiar el curso de la batalla.
    \item \textbf{Gestión de Recursos:} Los Puntos de Acción (PA) y Energía (PE) dictan el ritmo. Un guerrero agotado es un guerrero muerto.
    \item \textbf{Zonas Secas:} El uso de la magia tiene un costo físico en el campo de batalla, alterando el mapa de forma permanente.
\end{itemize}
\end{reglaespecial}

\section{Lo que Necesitas para Jugar}
Para comenzar tus batallas en la Era del Martillo, necesitarás lo siguiente:

\begin{description}
    \item[Miniaturas:] Que representen a tus héroes y escuadras.
    \item[Dados:] Un dado de doce caras (\textbf{1d12}) para pruebas de éxito, un dado de cien (\textbf{1d100}) para localización y un dado de seis (\textbf{1d6}) para el daño.
    \item[Superficie de Juego:] Un área de 90cm x 90cm con escenografía que represente ruinas, bosques o zonas de extracción de energía.
    \item[Cinta Métrica:] Las distancias se miden en centímetros (cm).
\end{description}

\section{El Concepto de Escaramuza}
A diferencia de las grandes batallas de bloques de infantería, este manual se enfoca en el **Combate Heroico**. Cada miniatura actúa de forma independiente. Una unidad puede ser un poderoso líder Evolgen o un pequeño grupo de apoyo Biskrorum, pero cada una tiene acceso a acciones, reacciones y habilidades únicas que consumen sus propios PA y PE.

\vfill
\begin{center}
    \textit{Prepárate, comandante. La Gema brilla, pero el terreno se seca.}
\end{center}
% --- ARCHIVO: Capitulos/sistema_activacion.tex ---

\chapter{Estructura del Turno}

El flujo de batalla en \textit{La Era del Martillo} se rige por el caos controlado de la guerra. No hay turnos de ejército completos; en su lugar, el combate es un toma y daca constante.

\section{Activación Alternada}
El juego se desarrolla a través de una serie de **Rondas**. En cada ronda, los jugadores alternan la activación de sus unidades.

\begin{reglaespecial}{El Estado de Agotamiento}
Al inicio de la Ronda, cada unidad recibe una \textbf{Ficha de Turno}. 
\begin{itemize}
    \item Cuando activas una unidad, esta consume su ficha.
    \item Una vez terminadas sus acciones, la unidad entra en estado \textbf{Agotado}.
    \item Una unidad Agotada no puede ser activada de nuevo hasta la siguiente ronda, a menos que una habilidad especial indique lo contrario.
\end{itemize}
\end{reglaespecial}

\section{Recursos de Unidad: PA y PE}
Cada miniatura tiene una reserva de energía y capacidad física que limita lo que puede hacer en un solo turno.

\subsection{Puntos de Acción (PA)}
Los **PA** representan el esfuerzo físico. Una unidad promedio comienza con su valor base de PA (indicado en su ficha) y puede gastarlos en acciones básicas o reacciones.

\subsection{Puntos de Energía (PE)}
Los **PE** representan el vínculo con la Gema o el esfuerzo mental. Se utilizan exclusivamente para habilidades mágicas o técnicas sobrehumanas. A diferencia de los PA, los PE pueden extraerse de tres fuentes:
\begin{enumerate}
    \item \textbf{Reserva Propia:} Energía almacenada en la unidad.
    \item \textbf{Vínculo de Gema:} Energía transmitida por un Líder (ver área de influencia).
    \item \textbf{El Entorno:} Extracción directa del mapa (creando una \textbf{Zona Seca}).
\end{enumerate}

\section{Acciones Disponibles}
Durante su activación, una unidad puede realizar cualquier combinación de acciones mientras tenga **PA** para costearlas.

\begin{table}[h]
\centering
\begin{tabularx}{\linewidth}{|l|c|X|}
\hline
\rowcolor{GoldRule!20} \textbf{Acción} & \textbf{Coste} & \textbf{Descripción} \\ \hline
\textbf{Mover} & 1 PA & Desplazarse hasta su valor de MOV en cm. \\ \hline
\textbf{Atacar} & 1 PA & Realizar un ataque de Melee o Rango. \\ \hline
\textbf{Placaje} & 1 PA & Intento de mover o derribar a un enemigo. \\ \hline
\textbf{Asistir} & 1 PA & Otorga +2 a la siguiente tirada de un aliado. \\ \hline
\textbf{Interactuar} & 1 PA & Robar la Gema, abrir puertas o usar objetos. \\ \hline
\end{tabularx}
\caption{Acciones Básicas de Combate}
\end{table}

\section{Reacciones}
Las unidades pueden actuar fuera de su activación gastando PA para responder a una amenaza enemiga (como realizar un Bloqueo o un Contraataque). Una unidad \textbf{Agotada} tiene limitado el uso de reacciones.

\vfill
\begin{center}
    \textit{La duda es el primer paso hacia la derrota. Gasta tus puntos con sabiduría.}
\end{center}


\begin{relato}{Crónica del Primer Impacto: Kaelen el Evolgen}
"Mis pulmones ardían, no por el aire viciado de las Tierras Secas, sino por la demanda de la Gema. Vi al Biskrorum alzando su maza; mi mente calculó mil trayectorias en un segundo. Gasté mi última chispa de energía en un parpadeo táctico. El acero chocó contra el vacío donde yo solía estar. No hay gloria en esto, solo la fría certeza de que Bhör nos mira, esperando que el último de nosotros caiga para declarar un ganador."
\end{relato}
% --- ARCHIVO: Capitulos/combate.tex ---

\chapter{El Arte de la Guerra}

El combate en \textit{La Era del Martillo} no es un simple intercambio de golpes; es un proceso de precisión donde la posición, la fuerza y la suerte dictan quién sobrevive.

\section{El Pool de Ataque}
Cuando una unidad declara una acción de **Melee** o **Rango**, el jugador lanza tres tipos de dados simultáneamente. Este conjunto de dados define la totalidad del ataque:

\begin{reglaespecial}{La Tirada de Triple Dado}
Al atacar, lanza estos dados juntos:
\begin{itemize}
    \item \textbf{1d12 (Éxito):} Se suma al atributo y al modificador que corresponda para superar la defensa del enemigo (\textbf{Esquiva} o \textbf{Bloqueo)}.
    \item \textbf{1d100 (Localización):} Determina en qué parte del cuerpo impacta el golpe.
    \item \textbf{1d6 (Daño):} Define la potencia base del impacto.
\end{itemize}
\end{reglaespecial}

\section{Resolución del Ataque}
Para determinar si un ataque impacta, se realiza una tirada enfrentada:

\begin{center}
    \textit{Ataque (1d12 + Atributo + Modificador) vs. Defensa (1d12 + Atributo + Modificador)}
\end{center}

Si el resultado del atacante es igual o superior al del defensor, el golpe conecta. En ataques de **Rango**, el defensor utiliza su atributo de \textbf{REF (Reflejos)}, mientras que en **Melee** suele utilizarse \textbf{BLOQ (Bloqueo)}.

\section{Localización y Daño}
Si el ataque impacta, consulta el resultado del **1d100** en la siguiente tabla para determinar las consecuencias físicas. Cada zona tiene un efecto si el daño supera la Resistencia de Daño (RD) del objetivo.

\begin{table}[h]
\centering
\begin{tabularx}{\linewidth}{|l|c|X|}
\hline
\rowcolor{WarhammerRed!20} \textbf{Zona} & \textbf{d100} & \textbf{Efecto Crítico} \\ \hline
Cabeza & 01-10 & Daño x2. Posible aturdimiento. \\ \hline
Torso  & 11-50 & Daño estándar. \\ \hline
Brazos & 51-75 & El objetivo no puede usar armas en ese brazo. \\ \hline
Piernas& 76-00 & El atributo de MOV se reduce a la mitad. \\ \hline
\end{tabularx}
\caption{Tabla de Localización de Daño}
\end{table}

\section{Resistencia y Reducción}
El daño final que recibe una unidad se calcula restando su estadística de **RD (Resistencia de Daño)** al valor obtenido en el dado de daño ($1d6$) más los bonificadores de arma.

\begin{tcolorbox}[colback=gray!10, colframe=black, title=Cálculo Final]
\centering
\textbf{Daño Resultante =\\ (1d6 + Bono Arma) - RD Enemiga}
\end{tcolorbox}

\begin{reglaespecial}{Incapacitación}
Si una unidad pierde todos sus puntos de vida en una extremidad específica, esa extremidad queda inutilizada permanentemente por el resto de la batalla, aplicando los penalizadores de la tabla de localización de forma constante.
\end{reglaespecial}
% --- ARCHIVO: Capitulos/magia_gema.tex ---

\chapter{La Gema y el Entorno}

La magia en \textit{La Era del Martillo} no es infinita. Es un recurso físico que se extrae de la tierra o se canaliza a través de artefactos antiguos. Comprender el flujo de la energía es la diferencia entre un místico poderoso y un cadáver agotado.

\section{Control de Gema (Gem Control)}
El **Control de Gema** representa la capacidad de un Líder o un portador para canalizar la energía de la Gema Primordial. 

\begin{reglaespecial}{Extracción de Energía}
Cada ronda, el portador de la Gema genera una cantidad de **PE (Puntos de Energía)** igual a su atributo de Control. 
\begin{itemize}
    \item Esta energía puede ser usada por el portador o distribuida a aliados.
    \item Los PE generados por la Gema no se acumulan entre rondas; lo que no se usa, se disipa al final de la activación del último aliado.
\end{itemize}
\end{reglaespecial}

\section{Área de Influencia}
El **Área de Influencia** es el radio táctico (medido en cm desde la peana del Líder) donde su autoridad y energía son tangibles.

\begin{itemize}
    \item \textbf{Canalización:} Las unidades dentro del área pueden usar los PE del Líder como si fueran propios.
    \item \textbf{Transferencia:} El Líder puede gastar una acción para "Cargar" a una unidad aliada. Esa unidad recibe PE que duran toda la ronda, permitiéndole salir del área de influencia y seguir usando habilidades mágicas sin coste adicional de terreno.
\end{itemize}

\section{El Costo del Mundo: Zonas Secas}
Cuando un hechicero o unidad necesita usar una habilidad y no tiene acceso a la energía de una Gema o reserva propia, debe recurrir a la **Extracción Forzada**.

\begin{reglaespecial}{Creación de una Zona Seca}
Si una unidad usa PE del entorno, el punto exacto donde se encuentra se convierte en una \textbf{Zona Seca}.
\begin{enumerate}
    \item Coloca un marcador de 3cm de radio centrado en la unidad actora.
    \item \textbf{Efecto:} Por el resto de la partida, ninguna unidad puede extraer energía del entorno dentro de ese radio.
    \item \textbf{Penalizador:} Las unidades que terminen su movimiento en una Zona Seca sufren un penalizador de -1 a sus tiradas de \textit{Focus}.
\end{enumerate}
\end{reglaespecial}

\section{Habilidades de Gema}
Las habilidades que consumen PE se dividen en:
\begin{description}
    \item[Hechizos Rápidos:] Se activan como Reacción (Coste en PA + PE).
    \item[Rituales:] Requieren una acción completa y suelen tener un alto consumo de PE o requieren sacrificar salud de la unidad.
    \item[Auras:] Habilidades pasivas que se mantienen activas mientras el Líder tenga al menos 1 PE en su reserva.
\end{description}

\vfill
\begin{center}
    \includegraphics[width=0.4\linewidth]{example-image-golden} % Aquí podrías poner un dibujo de la gema
    \\ \textit{"La tierra grita cuando la Gema calla."}
\end{center}

\end{multicols*}

\end{document}